\part{Critical evaluation, limitations and the hostile round}
\textit{Reconstructed section.}

\section{Critical evaluation}

\begin{qa}{C1}{What is the single biggest weakness of this thesis?}
No real-network evidence: everything is L3 in emulation or on a 5G core with a simulated RAN, and the traffic and
demands are modelled, not measured. Every quantitative claim is conditional on those models.
\end{qa}

\begin{qa}{C2}{What is the second biggest weakness?}
The headline mechanism comparison is null. The demand-aware allocator and RL are not significantly better than simple
slicing. The defensible contribution therefore shifts to ``service slicing with AQM, plus healing, measured
rigorously'', which is narrower than the original vision document claimed.
\end{qa}

\begin{qa}{C3}{Which original claims did you retract or narrow, and why?}
\begin{itemize}
\item ``AI-driven allocation improves service quality'': not supported (RL null, H3).
\item ``Context awareness drives the gains'': not supported (no measurable ablation effect).
\item ``RSNSD enables classification of rural traffic'': refuted for measured traffic (transfer F1 \NTrZero).
\item ``Self-healing with load balancing'': load balancing not implemented.
\item ``QoS Enforcer: prototype'' in the original status table: the Mininet path had never run when the verification
began.
\end{itemize}
\end{qa}

\begin{qa}{C4}{What result surprised you most?}
That HTB borrowing (a data-plane mechanism) erases most differences between allocation policies: once idle capacity is
lent at packet timescales, control-plane share computation matters only under strict caps or heavy overload. That
reframes where engineering effort should go: AQM and enforcement, not clever allocators.
\end{qa}

\begin{qa}{C5}{What would you do differently if you started again?}
\begin{itemize}
\item Build the measurement harness and baselines first, before mechanisms.
\item Run on real OVS from day one (the false-failover bug was invisible on bridges).
\item Pre-register hypotheses externally.
\item Seek real traces early.
\item Spend less on RL and RSNSD scale, more on real-world validity.
\end{itemize}
\end{qa}

\begin{qa}{C6}{Does the thesis generalise beyond the studied topology?}
Mechanism-level conclusions (slicing with AQM protects critical flows; loop-free protection switching gives sub-second
recovery) are standard enough to generalise to single-bottleneck networks. Quantitative results do not transfer to
multi-bottleneck, radio-limited or larger networks without new experiments.
\end{qa}

\begin{qa}{C7}{Is there enough here for a master's thesis?}
Yes, on method and evidence:
\begin{itemize}
\item a working integrated system;
\item a validated packet-level measurement method;
\item a 340-run, statistically analysed evaluation with negative results reported;
\item real-OVS and 5G-core replications;
\item documented discovery and correction of methodological errors.
\end{itemize}
The contribution is empirical and architectural rather than algorithmic, and it should be presented that way.
\end{qa}

\begin{qa}{C8}{What are the ethical implications of prioritising some services over others?}
Prioritisation encodes value judgements (whose traffic matters in a disaster) that should come from accountable local
stakeholders, not engineers. Floors guarantee minimum access for everyone; context weights should be governed,
transparent and auditable (the history log). Misuse (surveillance-driven classification, political prioritisation) is a
real risk. SliceNet classifies only coarse service classes from metadata.
\end{qa}

\begin{qa}{C9}{What are the privacy implications?}
Traffic-shape features are metadata but can still reveal activity. Deployment must capture headers only (or compute
features on the fly without storage), keep classification on site, and avoid per-user profiling. Our evaluation
captures included about 60\,B of synthetic payload for matching, which would be unacceptable on real traffic.
\end{qa}

\begin{qa}{C10}{What are the security risks?}
\begin{itemize}
\item Spoofed context alerts (forcing the disaster profile).
\item A compromised controller (arbitrary QoS, link shutdowns).
\item Evasion of SliceNet (mimicking healthcare traffic shapes to gain priority).
\end{itemize}
Mitigations: authenticated feeds, least privilege for controller commands, per-flow caps inside critical classes,
anomaly monitoring on classification.
\end{qa}

\begin{qa}{C11}{What would block a real rural deployment?}
Access to a standby backhaul (cost); skilled operators; power reliability for the controller and switches; integration
with the operator's 5G policy framework (PCF/QER instead of host tc); and stakeholder governance of priorities.
\end{qa}

\begin{qa}{C12}{What are the next six months of work?}
\begin{enumerate}
\item An L4 trial on an SDR-based private 5G testbed with a real backhaul bottleneck.
\item Real labelled traces and re-validation of SliceNet.
\item Standards-aligned enforcement (QER via PCF).
\item Controller resilience (state persistence, local fail-safe agent).
\item Multi-village topology with multipath healing (load balancing).
\end{enumerate}
\end{qa}

\begin{qa}{C13}{Which parts were AI-assisted, and how did you verify them?}
\textit{Answer truthfully for your own case.} Name the tools and say which components they helped with. The
verification was independent of authorship: unit tests on exact values, end-to-end runs on two switch implementations
and a 5G core, measurement cross-checks (floors, probe vs one-way delay), and the corrections documented in E12--E13.
Examiners care most that you understand and can defend every line. Rehearse the code walkthroughs in Q131--Q132.
\end{qa}

\section{The hostile round: one-line answers}

\begin{qa}{H1}{``Isn't this just Linux tc with extra steps?''}
The enforcement \emph{is} tc, deliberately. The thesis is about what to configure, when, and how to recover, and about
measuring that the outcome is right. The results say the configuring matters less than expected, and we report that.
\end{qa}

\begin{qa}{H2}{``Your RL doesn't work. Why is it in the thesis?''}
Because whether learning beats a good heuristic here is a legitimate research question, and a rigorously measured
negative answer (paired, Bonferroni-corrected, in emulation) is a result. It is presented as such, not as a
contribution.
\end{qa}

\begin{qa}{H3}{``Your classifier is trained on traffic you generated. Isn't the F1 meaningless?''}
It is an upper bound on separability of our models, evaluated on held-out runs and devices, with shortcuts removed.
The transfer study (\NTrZero{} zero-shot) shows exactly how generator-dependent such numbers are. Real-trace validation is
stated as the open gap.
\end{qa}

\begin{qa}{H4}{``Mininet isn't a real network.''}
Correct, so everything is labelled L3 in emulation. Mininet runs the real kernel stack, real TCP and real qdiscs; we
replicated the key results on real Open vSwitch and on a real 5G core; and we never claim field performance.
\end{qa}

\begin{qa}{H5}{``Most of your hypotheses failed. Is the thesis a failure?''}
No. Five pre-registered hypotheses were supported (ablation, healing outage, predictive switchover, SliceNet F1,
transfer gap). Six were not, and they are reported. A thesis that only confirms its hypotheses should worry a panel more.
\end{qa}

\begin{qa}{H6}{``Self-healing is just protection switching from the 1990s.''}
The switching is standard; what we add is coupling it to slice allocation (backup ceilings, degraded mode), loop-free
ordering without STP, OAM-gated fail-back, congestion-aware prediction, and packet-level outage measurement.
\end{qa}

\begin{qa}{H7}{``Your OAM check reads the emulator's netem configuration. Isn't that cheating?''}
It is a stand-in for Y.1731 or BFD measurement on the standby path, which reports the same delay and loss in reality.
It is documented as such in the code and the thesis. It only gates fail-back, and it does not affect any measured
outage.
\end{qa}

\begin{qa}{H8}{``How do I know you didn't tune on the test set?''}
No hyper-parameter search was done (fixed defaults). Thresholds were set by design and changed only for documented
bugs (which triggered full re-runs). SliceNet test groups are disjoint runs and devices.
\end{qa}

\begin{qa}{H9}{``Your 5G part isn't really 5G slicing.''}
Slice identity is genuine 5G (S-NSSAI, DNN, PDU sessions on Open5GS). Enforcement is CARS-SF on the UPF host, not
3GPP QoS. We say exactly that, and give the standards path (PCF, QER, RAN quotas).
\end{qa}

\begin{qa}{H10}{``If you had to cut half the thesis, what stays?''}
Service slicing with AQM and its measured critical-SLA benefit; loop-free self-healing with sub-second outage; the
measurement methodology. RL, RSNSD scale and the 5G stretch go.
\end{qa}

\begin{qa}{H11}{``What if I run it now and it fails?''}
\code{make setup data models test integration} runs without sudo. The real-OVS path runs in Docker
(\code{docker run --privileged cars-sf}). Every reported number comes from scripted pipelines over stored artefacts.
Offer a live run of Q528.
\end{qa}

\begin{qa}{H12}{``Why should a rural planner trust demands and SLAs you assumed?''}
They shouldn't trust the numbers; they can reuse the mechanism. Weights, floors and SLAs are registry data meant to be
set by local stakeholders, and the evaluation method can be re-run with their values.
\end{qa}

\begin{qa}{H13}{``Your sample sizes are tiny.''}
25 paired observations per main comparison, with CIs everywhere; 3-seed studies are presented descriptively. The null
allocator result is ``differences below about 0.02 in utility'', a bounded claim, not proof of zero.
\end{qa}

\begin{qa}{H14}{``Couldn't DiffServ marking with static queues give the same result?''}
Largely, for the steady state: our own static and 3GPP baselines are close. What it cannot do is healing coupled to
allocation, or slicing devices that cannot be trusted to mark (in-loop classification for the mixed UE).
\end{qa}
